\documentclass[a4paper, 10pt, twoside]{article}

\usepackage[top=1in, bottom=1in, left=1in, right=1in]{geometry}
\usepackage[utf8]{inputenc}
\usepackage[spanish, es-ucroman, es-noquoting]{babel}
\usepackage{setspace}
\usepackage{fancyhdr}
\usepackage{lastpage}
\usepackage{amsmath}
\usepackage{amsfonts}
\usepackage{amsthm}
\usepackage{verbatim}
\usepackage{graphicx}
\usepackage{float}
\usepackage{enumitem} % Provee macro \setlist
\usepackage{tabularx}
\usepackage{multirow}
\usepackage{hyperref}
\usepackage{multicol}
\usepackage[toc, page]{appendix}
\usepackage{ulem} % para subrayados especiales
\usepackage{color}
\usepackage{framed}
\usepackage{algorithm}
\usepackage{algpseudocode}
\usepackage{graphicx}
\usepackage{float}
\usepackage{listings}
\usepackage{hyperref}
%\DeclareUnicodeCharacter{00A0}{~}

\definecolor{shadecolor}{rgb}{0.95,0.95,0.95}



%%%%%%%%%% Configuración de Fancyhdr - Inicio %%%%%%%%%%
\pagestyle{fancy}
\thispagestyle{fancy}
\lhead{Trabajo Práctico 2 · Teoría de Lenguajes}
\rhead{Matayoshi · Panarello · Vega}
\renewcommand{\footrulewidth}{0.4pt}
\cfoot{\thepage /\pageref{LastPage}}

\fancypagestyle{caratula} {
   \fancyhf{}
   %\cfoot{\thepage /\pageref{LastPage}}
   \renewcommand{\headrulewidth}{0pt}
   \renewcommand{\footrulewidth}{0pt}
}
%%%%%%%%%% Configuración de Fancyhdr - Fin %%%%%%%%%%


%%%%%%%%%% Configuración de Hyperref - Inicio %%%%%%%%%%
\hypersetup{%
 % Para que el PDF se abra a pagina completa.
  pdfstartview= {FitH \hypercalcbp{\paperheight-\topmargin-1in-\headheight}},
  pdfauthor={Matayoshi,Panarello,Vega},
  pdfsubject={Tleng-TP1},
 %pdfkeywords={keyword1} {key2} {key3},
 colorlinks=true,
  linkcolor=black,
  urlcolor=blue
}
%%%%%%%%%% Configuración de Hyperref - Fin %%%%%%%%%%


%%%%%%%%%% Miscelánea - Inicio %%%%%%%%%%
% Evita que el documento se estire verticalmente para ocupar el espacio vacío
% en cada página.
\raggedbottom

% Deshabilita sangría en la primer línea de un párrafo.
\setlength{\parindent}{0em}

% Separación entre párrafos.
\setlength{\parskip}{0.5em}

% Separación entre elementos de listas.
\setlist{itemsep=0.5em}

% Asigna la traducción de la palabra 'Appendices'.
\renewcommand{\appendixtocname}{Apéndices}
\renewcommand{\appendixpagename}{Apéndices}
%%%%%%%%%% Miscelánea - Fin %%%%%%%%%%


%%%%%%%%%% Insertar diagrama - Inicio %%%%%%%%%%
\newcommand{\diagramav}[1]{
  \includegraphics[type=png,ext=.png,read=.png,width=16cm]{#1}
}

\newcommand{\diagramah}[1]{
  \includegraphics[type=png,ext=.png,read=.png,height=12.8cm,angle=90]{#1}
}
%%%%%%%%%% Insertar diagrama - Fin %%%%%%%%%%


\begin{document}
\pagenumbering{gobble}


%%%%%%%%%%%%%%%%%%%%%%%%%%%%%%%%%%%%%%%%%%%%%%%%%%%%%%%%%%%%%%%%%%%%%%%%%%%%%%%
%% Carátula                                                                  %%
%%%%%%%%%%%%%%%%%%%%%%%%%%%%%%%%%%%%%%%%%%%%%%%%%%%%%%%%%%%%%%%%%%%%%%%%%%%%%%%


\thispagestyle{caratula}

\begin{center}

\includegraphics[height=2cm]{caratula/DC.png}
\hfill
\includegraphics[height=2cm]{caratula/UBA.jpg}

\vspace{2cm}

Departamento de Computación,\\
Facultad de Ciencias Exactas y Naturales,\\
Universidad de Buenos Aires

\vspace{4cm}

\begin{Huge}
Trabajo Práctico 2
\end{Huge}

\vspace{0.5cm}

\begin{huge}
Compositor Musical
\end{huge}

\vspace{0.5cm}

\begin{Large}
Teoría de Lenguajes
\end{Large}

\vspace{1cm}

Primer Cuatrimestre de 2015

\vspace{3.5cm}

\begin{Large}
Grupo: Autores del Autómata Automático Autodestructivo
\end{Large}

\vspace{0.5cm}

\begin{tabular}{|c|c|c|}
\hline
Apellido y Nombre & LU & E-mail\\
\hline
Matayoshi, Leandro  & 79/11 & {\tt leandro.matayoshi@gmail.com}\\
Panarello, Bernabé  & 194/01 & {\tt bpanarello@gmail.com}\\
Vega, Leandro    & 698/11 & {\tt leandrogvega@gmail.com}\\
\hline
\end{tabular}

\end{center}

\newpage
\pagenumbering{arabic}

%%%%%%%%%%%%%%%%%%%%%%%%%%%%%%%%%%%%%%%%%%%%%%%%%%%%%%%%%%%%%%%%%%%%%%%%%%%%%%%
%% Índice                                                                    %%
%%%%%%%%%%%%%%%%%%%%%%%%%%%%%%%%%%%%%%%%%%%%%%%%%%%%%%%%%%%%%%%%%%%%%%%%%%%%%%%


\tableofcontents

\newpage


%%%%%%%%%%%%%%%%%%%%%%%%%%%%%%%%%%%%%%%%%%%%%%%%%%%%%%%%%%%%%%%%%%%%%%%%%%%%%%%
%% Descripción del informe			                                        %%
%%%%%%%%%%%%%%%%%%%%%%%%%%%%%%%%%%%%%%%%%%%%%%%%%%%%%%%%%%%%%%%%%%%%%%%%%%%%%%%

\section{Introducción del problema a resolver}
El objetivo de nuestro tp es, dado un archivo de entrada, parsearlo para poder tener un archivo de salida, respetando restricciones solicitadas para el correcto funcionamiento.

\subsection{Paso a paso introductivo de la resolución}
\begin{itemize}
\item En primera instancia definiremos una gramática necesaria para interpretar, de manera correcta, nuestro archivo de entrada.
\item Luego definiremos cada expresion regular como tokens en el lexer\_rules, de esta manera cargamos en la computadora nuestras expresiones regulares definidas en la gramática.
\item El siguiente paso será definir las producciones correspondientes, las cuales gracias a los tokens definidos en el paso anterior, podremos construir, diferenciarlas una de las otras y filtrar aquellas que no sean válidas.
\item Cada producción llamará a su función asociada, formando de esta manera un árbol AST, desde las hojas hasta su raiz, usando en cada una atributos sintetizados para poder intercambiar valores de una rama a la otra y poder validar las condiciones especificadas en nuestro lenguaje y el enunciado.
\item Una vez corroborado y habiendo obtenido con éxito todo lo anterior, procederemos a escribir el archivo MIDI de acuerdo al formato del mismo sin dificultades.
\end{itemize}

\newpage

\section{Descripción del problema resuelto}
En esta sección explicaremos cada parte implementada para realizar los procedimientos requeridos. Contaremos dudas, errores que fueron surgiendo y explicaremos las decisiones tomadas. Para eso vamos a dividirlo en cuatro secciones que detallamos a continuación ordenadas de las formas en las que lo fuimos realizando.

\subsection{Gramática}
No contamos con demasiadas dificultades, miramos cada paso de la descripción de la partitura y fuimos creando las producciones necesarias. Se explicita en la sección Gramática.

\subsection{PLY}
En nuestra implementación utilizamos las herramientas brindadas por PLY como explicamos en la introducción, estas son:
\subsubsection{Lexer}
Vamos a utilizar un lexer para poder tener definidas nuestras expresiones regulares y poder decidir que cadenas son o no válidas según nuestra gramática.
Para generar un lexer vamos a realizar lo descripto en el siguiente link \url{www.dc.uba.ar/materias/tl/2015/c1/files/tp2-clase-intro-a-ply/at_download/file}.
Crearemos un archivo lexer\_rules.py, definiendo los tokens y las reglas.
La lista de tokens será explicitada en la sección de la gramática.\newline
Luego pasamos a definir las reglas para cada expresión regular que deseamos caputar. En esta parte tuvimos bastantes problemas, principalmente porque optamos por usar reglas 'simples' en todas las reglas.\newline
El primero de los problemas fue al hacer la siguiente expresión regular "$(do|re|mi|fa|sol|la|si)(+|-)?$". Al tenerla toda junta, la regla matcheaba con todas las notas hasta el final del primer paréntesis sin mirar lo que continuaba y, en caso de tener en el archivo de entrada "sol+" en una nota, buscara una regla que inicialice con un + o - y no la encontrará lógicamente. La solución fue separarlas, pudiendo solucionar el problema mencionado. \newline
El segundo problema que surgió fue el tema del orden, al tener reglas como const que generan todo el alfabeto, nos matcheaba voz, compas, entre otras, cuando nosotros en realidad queriamos que esas palabras matcheen en otra regla definida. Como nosotros teníamos definido reglas 'simples', al tratar de ordenar las reglas notamos que la lógica de la clase re (regular expression) tomaba como primera a la que definía en su regla el string más largo, esto nos hizo rever la forma de definir las reglas. Mirando y testeando la clase re, pudimos corroborar que, definiendo las reglas como funciones, se respetaba el orden en las que eran definidas en el script. De esta manera pasamos todas las reglas a funciones como se explica en el pdf del link para reglas 'complejas', logrando salvar dicho problema.
\subsubsection{Parser}
Vamos a utilizar un parser para poder darle una estructura y una semántica a nuestra grámatica.
Para generar un parser vamos a realizar lo descripto en el link mencionado en el lexer desde la página 15 en adelante. Creamos un archivo parser\_rules.py, en el cual vamos a definir las producciones de la gramática y cómo generar el arbol AST (abstract Syntax Tree). Este archivo usará los tokens de lexer\_rules para construir las producciones, diferenciarlas una de las otras y filtrar aquellas que no sean válidas.\newline
Cada producción llamará a su función interna, que estará definida en parserobject.py, formando el árbol AST desde las hojas hasta su raiz. Cada una de ellas creará un objeto nodo y usará atributos sintetizados para poder intercambiar valores de una rama a la otra y poder validar las condiciones especificadas en nuestro lenguaje, en otras palabras le estaremos dando una semántica a nuestras producciones.
\newpage

\section{Gramática}
%\subsection{Gramática derivada del enunciado}
%
%H $\rightarrow$ Tempo Compasheader (Constinit)? Voicelist\newline
%
%Tempo $\rightarrow$ \# tempo Shape Num\newline
%
%Compasheader $\rightarrow$ \# compas Num / Num\newline
%
%Constinit $\rightarrow$ Constlist\newline
%
%Constlist $\rightarrow$ (Const)$^+$\newline
%
%Const $\rightarrow$ const Value = Value ;\newline
%
%Voicelist $\rightarrow$ (Voice)$^+$\newline
%
%Voice $\rightarrow$ voz ( Value ) \{ VoiceContent \}\newline
%
%Voicecontent $\rightarrow$ (Compas$|$Compasloop)$^+$\newline
%
%Compasloop $\rightarrow$ repetir ( Value ) \{ Compaslist \}\newline
%
%Compas $\rightarrow$ compas \{ Notelist \}\newline
%
%Notelist $\rightarrow$ (Note$|$Silence)$^+$\newline
%
%Note $\rightarrow$ nota ( Notename (+$|$-)? , Value , Shape (.)? ) ;\newline
%
%Value $\rightarrow$ Cname$|$Num\newline
%
%Silence $\rightarrow$ silencio ( Shape (.)? ) ;\newline
%
%Shape $\rightarrow$ blanca$|$negra$|$redonda$|$semicorchea$|$corchea$|$semifusa$|$fusa \newline
%
%Notename $\rightarrow$ do$|$re$|$mi$|$fa$|$sol$|$la$|$si \newline
%
%Cname $\rightarrow$ ([a-z] $|$ [A-Z]) ([0-9] $|$ [a-z] $|$ [A-Z])$^*$\newline
%
%Num $\rightarrow$ [0] $|$ ([1-9] [0-9])$^*$\newline

\subsection{Gramática deducida}
H $\rightarrow$ \{TEMPO\}\{COMPASHEADER\}\{CONSTINIT\}\{VOICELIST\}\newline
H $\rightarrow$ \{TEMPO\}\{COMPASHEADER\}\{VOICELIST\}\newline

TEMPO $\rightarrow$ \{tempobegin\}\{shape\}\{num\}\newline

COMPASHEADER $\rightarrow$ \{compasheaderbegin\}\{num\}\{slash\}\{num\}\newline

CONSTINIT $\rightarrow$ \{CONSTLIST\}\newline

CONSTLIST $\rightarrow$ \{CONST\}\newline
CONSTLIST $\rightarrow$ \{CONSTLIST\}\{CONST\}\newline

CONST $\rightarrow$ \{const\}\{VALUE\}\{equals\}\{VALUE\}\{semicolon\}\newline

VOICELIST $\rightarrow$ \{VOICE\}\newline
VOICELIST $\rightarrow$ \{VOICELIST\}\{VOICE\}\newline

VOICE $\rightarrow$ \{voicebegin\}\{leftpar\}\{VALUE\} \{rightpar\}\{leftcurl\}\{VOICECONTENT\}\{rightcurl\}\newline

VOICECONTENT $\rightarrow$ \{COMPAS\}\newline
VOICECONTENT $\rightarrow$ \{COMPASLOOP\}\newline
VOICECONTENT $\rightarrow$ \{VOICECONTENT\}\{COMPAS\}\newline
VOICECONTENT $\rightarrow$ \{VOICECONTENT\}\{COMPASLOOP\}\newline

COMPASLOOP $\rightarrow$ \{loopbegin\}\{leftpar\}\{VALUE\}\ {rightpar\}\{leftcurl\}\{COMPASLIST\}\{rightcurl\}\newline

COMPASLIST $\rightarrow$ \{COMPAS\}\newline
COMPASLIST $\rightarrow$ \{COMPASLIST\}\{COMPAS\}\newline

COMPAS $\rightarrow$ \{compasbegin\}\{leftcurl\}\{NOTELIST\}\{rightcurl\}\newline

NOTELIST $\rightarrow$ \{NOTE\}\newline
NOTELIST $\rightarrow$ \{SILENCE\}\newline
NOTELIST $\rightarrow$ \{NOTELIST\}\{NOTE\}\newline
NOTELIST $\rightarrow$ \{NOTELIST\}\{SILENCE\}\newline

NOTE $\rightarrow$ \{notebegin\}\{leftpar\}\{notename\}\{alter\}\{comma\}\{VALUE\} \{comma\}\{shape\}\{puntillo\}\{rightpar\}\{semicolon\}\newline
NOTE $\rightarrow$ \{notebegin\}\{leftpar\}\{notename\}\{alter\}\{comma\}\{VALUE\} \{comma\}\{shape\}\{rightpar\}\{semicolon\}\newline
NOTE $\rightarrow$ \{notebegin\}\{leftpar\}\{notename\}\{comma\}\{VALUE\} \{comma\}\{shape\}\{puntillo\}\{rightpar\}\{semicolon\}\newline
NOTE $\rightarrow$ \{notebegin\}\{leftpar\}\{notename\}\{comma\}\{VALUE\} \{comma\}\{shape\}\{rightpar\}\{semicolon\}\newline

SILENCE $\rightarrow$ \{silencebegin\}\{leftpar\}\{shape\}\{rightpar\}\{semicolon\}\newline
SILENCE $\rightarrow$ \{silencebegin\}\{leftpar\}\{shape\}\{puntillo\}\{rightpar\}\{semicolon\}\newline

VALUE $\rightarrow$ \{cname\}\newline
VALUE $\rightarrow$ \{num\}\newline

\subsection{Tokens}
\begin{itemize}
\item TEMPOBEGIN: \# tempo

\item CONST: const

\item EQUALS: $=$

\item SEMICOLON: ;

\item VOICEBEGIN: voz

\item LEFTPAR: (

\item RIGHTPAR: )

\item LEFTCURL: \{

\item RIGHTCURL: \}

\item COMPASHEADERBEGIN: \# compas

\item COMPASBEGIN: compas

\item LOOPBEGIN: repetir

\item SLASH: $/$

\item NOTEBEGIN: nota

\item SILENCEBEGIN: silencio

\item PUNTO: .

\item ALTER: $+|-$

\item SHAPE: $blanca|negra|redonda|semicorchea|corchea|semifusa|fusa$

\item NOTENAME: $do|re|mi|fa|sol|la|si$

\item COMMA: ,

\item CNAME: $(([a-z]|[A-Z])([0-9]|[a-z]|[A-Z])*)$

\item NUM: $([0]|[1-9][0-9]*)$

\end{itemize}

\newpage

\section{Tests}
\lstset{
  frame=none,
  xleftmargin=2pt,
  stepnumber=1,
  numbers=left,
  numbersep=5pt,
  numberstyle=\ttfamily\tiny\color[gray]{0.3},
  belowcaptionskip=\bigskipamount,
  captionpos=b,
  escapeinside={*'}{'*},
  language=prolog,
  tabsize=2,
  emphstyle={\bf},
  commentstyle=\it,
  stringstyle=\mdseries\rmfamily,
  showspaces=false,
  keywordstyle=\bfseries\rmfamily,
  columns=flexible,
  basicstyle=\small\sffamily,
  showstringspaces=false,
  morecomment=[l]\%,
}

\subsection{Tests con fallas}

\subsubsection{Test 1: Tiene compaces con distinta duración}

\begin{lstlisting}
#tempo negra 30
#compas 3/4

const oct1 = 2;
const oct2 = 6;
const oct3 = 1;

// Instrumentos
const flauta = 51;

voz (flauta)
{
  compas
  {
    nota(do, oct3, blanca.);
    nota(re, oct1, redonda);
  }

  compas
  {
    nota(mi, oct2, blanca);
    nota(la, oct1, negra);
  }
}
\end{lstlisting}
\vspace{5mm}

\subsubsection{Test 2: Tiene voces con compaces de distinta duración}

\begin{lstlisting}
#tempo negra 120
#compas 2/2

const oct1 = 5;
const oct2 = 2;
const oct3 = 4;

// Instrumentos
const violin = 20;
const guitarra = 12;

voz (violin)
{
  compas
  {
    nota(do, oct3, blanca.);
    nota(re, oct1, negra);
  }

  compas
  {
    nota(mi, oct2, blanca.);
    nota(la, oct1, negra);
  }
}

voz (guitarra)
{
  compas
  {
    nota(do, oct3, fusa);
    nota(re, oct1, semifusa.);
  }

  compas
  {
    nota(mi, oct2, fusa);
    nota(la, oct1, semifusa.);
  }
}
\end{lstlisting}
\vspace{5mm}

\subsubsection{Test 3: Constante que apunta a una constante no definida}

\begin{lstlisting}
#tempo negra 60
#compas 1/1

const oct1 = 3;
const oct2 = ConstanteTrucha;

// Instrumentos
const bajo = 20;

voz (bajo)
{
  compas
  {
    nota(do, oct1, blanca.);
    nota(re, oct1, negra);
  }

  compas
  {
    nota(mi, oct2, blanca.);
    nota(la, oct2, negra);
  }
}
\end{lstlisting}
\vspace{5mm}

\subsubsection{Test 4: Constante definida circularmente}

\begin{lstlisting}
#tempo negra 60
#compas 2/8

const oct1 = 3;
const oct2 = 5;

// Instrumentos
const bajo = 20;
const malPensado = malPensado;

voz (bajo)
{
  compas
  {
    nota(do, oct1, corchea.);
    nota(re, oct1, semicorchea);
  }

  compas
  {
    nota(mi, oct2, semicorchea);
    nota(la, oct2, corchea.);
  }
}
\end{lstlisting}
\vspace{5mm}

\subsection{Tests correctos}

\subsubsection{Test 1: Simple}

\begin{lstlisting}
#tempo negra 30
#compas 2/4

const oct1 = 2;
const oct2 = 6;
const oct3 = 1;

// Instrumentos
const flauta = 51;

voz (flauta)
{
  compas
  {
    nota(do, oct3, negra);
    nota(re, oct1, negra);
  }

  compas
  {
    nota(mi, oct2, blanca);
  }
}
\end{lstlisting}
\vspace{5mm}

\subsubsection{Test 2: Con varias voces}

\begin{lstlisting}
#tempo negra 120
#compas 3/4

const oct1 = 2;
const oct2 = 6;
const oct3 = 1;
const oct4 = 3;

// Instrumentos
const piano = 65;
const violin = 31;

voz (piano)
{
  compas
  {
    nota(sol, oct3, blanca);
    nota(fa+, oct2, negra);
  }

  compas
  {
    nota(mi, oct2, negra.);
    nota(fa, oct4, corchea);
    nota(mi, oct2, negra);
  }
  compas
  {
    silencio(negra);
    nota(sol-, oct1, negra);
    nota(sol-, oct1, negra);
  }
}

voz (violin)
{
  compas
  {
    nota(la, oct1, blanca.);
  }
  
  compas
  {
    nota(mi, oct2, negra.);
    nota(fa, oct4, corchea);
    nota(mi, oct2, negra);
  }
  compas
  {
    silencio(semicorchea);
    nota(mi, oct1, semicorchea);
    nota(sol, oct4, corchea);
    nota(sol, oct4, blanca);
  }
}
\end{lstlisting}
\vspace{5mm}

\subsubsection{Test 3: Con repeticiones}

\begin{lstlisting}
#tempo negra 120
#compas 2/4

const oct1 = 2;
const oct2 = 6;
const oct3 = 1;
const oct4 = 3;

// Instrumentos
const flauta = 10;
const violin = 31;

voz (flauta)
{
  repetir (5)
  {
  	compas
  	{
    		nota(sol, oct3, negra);
    		nota(fa+, oct2, negra);
  	}

  	compas
  	{
    		nota(mi, oct2, negra.);
    		nota(fa, oct4, corchea);
  	}
  	compas
  	{
    		nota(sol, oct1, blanca);
  	}
  }
}

voz (violin)
{
  compas
  {
    nota(la, oct1, blanca);
  }
  
  compas
  {
    nota(fa, oct4, corchea);
    nota(mi, oct2, negra.);
  }
  compas
  {
    silencio(semifusa);
    nota(si+, oct1, semifusa);
    nota(si, oct4, semifusa);
    nota(fa, oct4, semifusa);
    nota(sol, oct4, semicorchea);
    nota(mi-, oct4, corchea);
    nota(re, oct4, negra);
  }
}
\end{lstlisting}
\vspace{5mm}
\newpage

\section{Manual del programa}
\subsection{Modo de uso}
Línea de ejecución: ./musileng entrada.mus salida.txt

\subsubsection{Reglas para evitar posibles errores}
Se aceptan archivos de entrada que contengan:
\begin{itemize}
\item Todos los compases deben tener la misma duración al sumar la duración de sus notas y/o silencios.
\item Todas las voces deben tener la misma cantidad de compases.
\item No deben existir constantes indefinidas o definidas circularmente. Ejemplo constante no definida: $const$ $eval = hola;$ (hola jamás se definió). Ejemplo constante definida circularmente: $const$ $eval1 = eval2$ $;$ $const$ $eval2 = eval1$ $;$.
\item La suma de la duración de cada compas debe ser igual a $num1/num2$, donde $num1$ y $num2$ son los números definidos en $\#compas$ $num1/num2$ en el encabezado.
\item Una constante definida como instrumento sólo acepta valores del $1$ al $127$.
\item El valor colocado en 'repetir' debe ser mayor a $0$.
\item Una constante definida como octava sólo acepta valores del $1$ al $9$.
\item No debe haber una constante definida más de una vez.
\item En $\#compas$ $num1/num2$ del encabezado, $num2$ debe ser un número que corresponda al de las figuras definidas, para corroborar puede ver la página 2 de \url{www.dc.uba.ar/materias/tl/2015/c1/tp2-enunciado-compositor-musical/at_download/file} en la tabla Figura y Valor numérico asociado.
\end{itemize}
En caso contrario que no se respete lo mencionado anteriormente, nuestro programa especificará el error cometido para que pueda solucionarlo.

Para más información sobre los archivos de entrada y salida puede mirarse el siguiente pdf \url{www.dc.uba.ar/materias/tl/2015/c1/tp2-enunciado-compositor-musical/at_download/file}.

\subsection{Requerimientos necesarios para ejecutar}
\begin{itemize}
\item Programa: Python
\item Versión: 2.7
\end{itemize}

%\begin{itemize}
%\item En la primera línea: #tempo FIGURA N (Falta desarrollar)
%\item En la segunda línea: #compas N/M
%\item En las siguientes líneas definiremos las constantes de las octavas, por ejemplo: "$const oct1 = 6;$"
%\item Luego de las octavas definiremos las constantes de instrumentos, por ejemplo: "$const piano = 31;$"
%\item Más adelante seguirán una lista de voces, a las cuales se les asignarán su respectivo instrumento definido anteriormente
%\item En cada voz definiremos una lista de compases o un loop de compases.
%\item Un loop de compases se define mediante un repetir, el cual cuenta con un número entero que indicará la cantidad de veces que se repetirá, seguido de una lista de compases que son los que se repetirán la cantidad de mencionada en el número entero.
%\item Un compás tiene una lista de notas o/y silencios.
%\item 
%\end{itemize}
\newpage

\section{Conclusiones}
\begin{itemize}
\item PLY es una herramienta muy útil, facilita el parseo y nos permite, de manera mucho más corta y sencilla, realizar las diferentes tareas para nuestro lenguaje.
\item Lo más complicado fue la parte del lexer, específicamente en la parte que definimos los tokens. Porque como explicamos en la descripción necesitabamos establecer un orden y para eso tuvimos que, entre otras cosas, entender como funcionaba y probar la clase re (regular expression) de python.
\item Los temas aportados en las clases, como gramática de atributos, TDS, parser, gramática LALR fueron útiles para poder resolver los problemas presentados.

\end{itemize}
\newpage

\section{Apéndices}

\subsection{musileng.py}
\lstset{
 numbers=left,
 numbersep=5pt,
 numberstyle=\ttfamily\tiny\color[gray]{0.3},
 basicstyle=\ttfamily,
 columns=fullflexible,
 keepspaces=true,
 breaklines=true,
}

\begin{verbatim}
#!/usr/bin/env python
from parserobjects import *


#Programa principal MUSILEN

import lexer_rules
import parser_rules

from sys import argv, exit
from midi import *
from ply.lex import lex
from ply.yacc import yacc


if __name__ == "__main__":
    if len(argv) != 3:
        print ("Invalid arguments.")
        print ("Use:")
        print ("  ./musileng.py <archivo_entrada> <archivo_salida>")
        exit(1)

    try:
        inputfile = argv[1]
        outputfile = argv[2]

        with open (inputfile, "r") as myfile:
            text=myfile.read()

        ConstantManager.createInstance([],[])

        lexer = lex(module=lexer_rules)
        parser = yacc(module=parser_rules)

        expression = parser.parse(text, lexer)
        midi = MidiTranslator()

        with open (outputfile, "w") as file_output:
            midi.generateMIDIFile(expression,file_output)


    except Exception as ex:
        print('ERROR:')
        for v in ex.args:
            print(str(v))

        exit(1)
\end{verbatim}


\newpage

\subsection{lexer\_rules.py}

\begin{verbatim}
tokens = [
   'COMMENT',
   'NUM',
   'CONST',
   'EQUALS',
   'CNAME',
   'SILENCEBEGIN',
   'SEMICOLON',
   'VOICEBEGIN',
   'LEFTPAR',
   'RIGHTPAR',
   'LEFTCURL',
   'RIGHTCURL',
   'COMPASBEGIN',
   'COMPASHEADERBEGIN',
   'LOOPBEGIN',
   'SLASH',
   'NOTEBEGIN',
   'SHAPE',
   'PUNTO',
   'ALTER',
   'NOTENAME',
   'COMMA',
   'NEWLINE',
   'TEMPOBEGIN'
]

#Ignoro espacios y tabs
t_ignore  = ' \t'

#Ignorar comoentarios
def t_COMMENT(token):
  r'//.*'

def t_NEWLINE(token):
  r"\n+"
  token.lexer.lineno += len(token.value)

#Chequear regexs con: https://regex101.com/#python.
# Tener en cuenta poner las expresiones que matchean 
# strings mas largos primero cuando hay ambiguedad.
# Usamos '\s' para matchear espacios en blanco.
def t_TEMPOBEGIN(token):
  r"\#tempo"
  return token

def t_CONST(token):
  r"const"
  return token

def t_EQUALS(token):
  r"\="
  return token

def t_SEMICOLON(token):
  r"\;"
  #print ("semicolon")
  return token

def t_VOICEBEGIN(token):
  r"voz"
  return token

def t_LEFTPAR(token):
  r"\("
  return token

def t_RIGHTPAR(token):
  r"\)"
  return token

def t_LEFTCURL(token):
  r"\{"
  return token

def t_RIGHTCURL(token):
  r"\}"
  return token

def t_COMPASHEADERBEGIN(token):
  r"\#compas"
  return token

def t_COMPASBEGIN(token):
  r"compas"
  return token

def t_LOOPBEGIN(token):
  r"repetir"
  return token

def t_SLASH(token):
  r"/"
  return token

def t_NOTEBEGIN(token):
  r"nota"
  return token

def t_SILENCEBEGIN(token):
  r"silencio"
  return token

def t_PUNTO(token):
  r"\."
  return token

def t_ALTER(token):
  r"\+|\-"
  return token

def t_SHAPE(token):
  r"blanca|negra|redonda|semicorchea|corchea|semifusa|fusa"
  return token

def t_NOTENAME(token):
  r"do|re|mi|fa|sol|la|si"
  return token

def t_COMMA(token):
  r"\,"
  return token

def t_CNAME(token):
  r"(([a-z]|[A-Z])([0-9]|[a-z]|[A-Z])*)"
  return token

def t_NUM(token):
    r"([0]|[1-9][0-9]*)"
    token.value = int(token.value)
    return token

def t_error(token):
    message = "[Lexer] Token desconocido"
    message += "\nline:" + str(token.lineno)
    message += "\nposition:" + str(token.lexpos)
    raise Exception(message)
\end{verbatim}
\newpage

\subsection{parserobjects.py}
\begin{verbatim}
#!/usr/bin/python
# -*- coding: latin-1 -*-

#En este archivo implementamos todas las clases 
#que representan los objetos que se van parseando (no terminales),
#algunas clases auxiliares y el manejador de constantes.

#Clase abstracta que representa expresiones (nodos en a gramatica)
class Expression(object):
    pass

#Representa una lista como un valor y un puntero a otra lista
#Por ejemplo, A -> id, A -> A id.
# Notar que los elementos de la lista estan en orden inverso
class ExpressionList(Expression):
    def __init__(self, current, nextList):
        self._list = nextList + [current]

    def getList(self):
        return self._list

#Representa un nodo del elemento const en la gramatica
class Const(Expression):
    def __init__(self, cname, val, isPointer):
        self._cname = cname
        self._value = val
        self._isPointer = isPointer

    def getName(self):
        return self._cname

    def getValue(self):
        return self._value

    #indica si la constante apunta a otra constante
    def isPointer(self):
        return self._isPointer


#Representa una lista de constantes
class ConstList(ExpressionList):
    pass

#Representa la lista de contenido de una voz (compaces o loops)
class VoiceContent(ExpressionList):
    def __init__(self, current, nextList):
        if nextList != [] and current.getDuration() != nextList[0].getDuration():
            raise Exception("Se detecto un compas dentro de una voz con duracion
           	               distinta al resto de los compases dentro de la misma.")

        self._duration = current.getDuration()
        super(VoiceContent,self).__init__(current, nextList)

    def getDuration(self):
        return self._duration

#Representa una lista de constantes dentro de un loop
class CompasList(ExpressionList):
    def __init__(self, current, nextList):
        if nextList != [] and current.getDuration() != nextList[0].getDuration():
            raise Exception("Se detectaron compases con distinta duracion dentro 
                           de un bucle.")

        self._duration = current.getDuration()
        super(CompasList,self).__init__(current, nextList)

    def getDuration(self):
        return self._duration

#Manejador de constantes. Mantiene una lista de todas las constantes declaradas 
#y permite obtener
#los valores de las mismas dado un nombre de constante.
class ConstantManager:
    #Inicializa el manejador de constantes dada una lista de objetos tipo Const.

    @staticmethod
    def createInstance(constList, reserved):
        ConstantManager._instance = ConstantManager(constList, reserved)

    @staticmethod
    def getInstance():
        if ConstantManager._instance == None:
            ConstantManager.createInstance([],[])
        return ConstantManager._instance


    #Inicializacion y validacion de consistencia de la declaracion de constantes
    def __init__(self, constList, reserved):
        self.dictConst = {}
        for c in constList:
            if reserved.count(c.getName()) > 0:
                raise Exception("El nombre de constante <{0}> es una palabra
                               reservada".format(c.getName()))

            self.Add (c.getName(), c)

        #Pido los valores de todas las ctes para que se validen
        for c in constList:
            self.getValue(c.getName())

    def Add(self, cname, const):
        if self.dictConst.get(cname) != None:
            raise Exception("Constante duplicada: {0}".format(cname))

        self.dictConst[cname] = const

    #Devuelve el valor de una constante. 
    #Si es una referencia a otra constante, calcula primero el valor recursivamente
    def getValue(self, cname):
        const = self.dictConst.get(cname)
        if const == None:
            raise Exception("Constante no definida: {0}".format(cname))
        else:
            constVal = self.resolveVal(const)

        return constVal


    #Calcula el valor de una constante. Si es numerica (caso base), devuelve su valor.
    #Si es un puntero, busca el valor de la constante a la que apunta (recursivamente), 
    #chequeando que no hayan loops.
    def resolveVal(self, const):
        if (not const._isPointer):
            return const.getValue()

        found = False
        visited = [const.getName()]
        next = const
        while(not found):
            if (next.isPointer()):
                nextcname = next.getValue()
                if visited.count(nextcname) > 0:
                    raise Exception("Definicion circular de constante
                                    <{0}>".format(const.getName()))

                next = self.dictConst.get(nextcname)
                if (next == None):
                    raise Exception("La constante <{0}> apunta (directa
                                   o indirectamente) al nombre <{1}> que 
                                   no esta definido".format(const.getName(),nextcname))
                visited.append(nextcname)
            else:
                #const.setValue(next.getValue())
                found = True

        return next.getValue()



class Tempo(Expression):
    def __init__(self, shape, num):
        self._shape = shape
        self._num = num

    def getFigure(self):
        return self._shape

    def getCount(self):
        return self._num

    def getShapeDuration(self):
        return duration_from_shape(self.getFigure())

#Representa el nodo que define la duracion de los compases para toda la partitura
class CompasHeader(Expression):
    def __init__(self, num1, num2):
        if num2 == 0:
            raise Exception("El denominador del compas debe ser mayor a cero")
        if not num2 in (1,2,4,8,16,32,64):
            raise Exception("El denominador definido para el tempo de los compaces
                            no corresponde a una figura valida.")


        self._numerator = num1
        self._denominator = num2


    #Devuelve la duracion del compas, en redondas
    def getDuration(self):
        return float(self._numerator) / float(self._denominator)


    def getDenominator(self):
        return self._denominator


    def getNumerator(self):
        return self._numerator

#Representa un nodo voz(N)
class Voice(Expression):
    def __init__(self, value, voiceContent):
        if value > 127:
          raise Exception("El instrumento {0} es invalido".format(value))

        self._value = value
        self._compasses = voiceContent.getList()

    def getInstrument(self):
        return self._value

    def getCompasses(self):
        return self._compasses


#Representa un nodo repetir(N) con una lista de compaces adentro
class CompasLoop(Expression):
    def __init__(self, value, compasList):
        if value < 1:
            raise Exception("Existe un bucle con cantidad de repeticiones 0. 
                            Los bucles deben tener al menos una repeticion.")

        self._compasList = compasList
        self._value = value
        self._duration = compasList.getDuration()


    def getRepeat(self):
        return self._value

    def getCompasses(self):
        return self._compasList

    def getDuration(self):
        return self._duration

    def isLoop(self):
        return True


def duration_from_shape(shape):
    if shape == 'redonda':
        duration = 1.0
    elif shape == 'blanca':
        duration = 1.0 / 2.0
    elif shape == 'negra':
        duration = 1.0 / 4.0
    elif shape == 'corchea':
        duration = 1.0 / 8.0
    elif shape == 'semicorchea':
        duration = 1.0 / 16.0
    elif shape == 'fusa':
        duration = 1.0 / 32.0
    elif shape == 'semifusa':
        duration = 1.0 / 64.0
    else:
        raise Exception("Figura no definida")
    return duration


class Silence(Expression):
    def __init__(self, duration, puntillo):
        self._duration = duration_from_shape(duration)
        if puntillo:
            self._duration = self._duration * (3.0 / 2.0)

    def getDuration(self):
        return self._duration

    def getShape(self):
        return self._shape

    def isSilence(self):
        return True


class Note(Expression):
    def __init__(self, notename, alter, value, shape, puntillo):
        if value < 1 or value > 9:
            raise Exception("La octava {0} es invalida".format(value))

        # concatenacion de string
        if alter != None:
            self._height = notename + alter
        else:
            self._height = notename

        self._value = value
        self._duration = duration_from_shape(shape)

        if puntillo:
            self._duration = self._duration * (3.0 / 2.0)


    def getOctave(self):
        return self._value

    def getHeight(self):
        return self._height

    def getDuration(self):
        return self._duration

    def isSilence(self):
        return False

#Objeto nodo principal
class Root(Expression):
    def __init__(self, tempo, compasheader,constlist,voicelist):
        ### El tiempo del primer compas de la primera voz debe coincidir 
        ### con el compasheader
        if voicelist.getList()[0].getCompasses()[0].getDuration()
        != compasheader.getDuration():
          raise Exception("La duracion de los compases es distinta
                          a la del encabezado")
        self._tempo = tempo
        self._compasheader = compasheader
        self._constlist = constlist
        self._voicelist = voicelist

    def getDuration(self):
        return self._tempo

    def getTempo(self):
        return self._tempo

    def getCompasHeader(self):
        return self._compasheader

    def getConstList(self):
        return self._constlist

    def getVoiceList(self):
        return self._voicelist

class VoiceList(Expression):
    def __init__(self, voice, voicelist=[]):
        if len(voicelist)>0:
            if (voice.getCompasses()[0].getDuration() !=
            voicelist[0].getCompasses()[0].getDuration()):
                raise Exception("La duracion de los compases de las voces 
                                son distintas")
            if len(voicelist)>=16:
            	raise Exception("La cantidad de voces supera las 16")
        self._voicelist = voicelist + [voice]

    def getList(self):
        return self._voicelist

class Compas(Expression):
    def __init__(self, notelist):
        self._notelist = notelist
        ### Calculo la duracion del compas
        duration = 0.0
        for note in notelist:
            duration += note.getDuration()
        self._duration = duration

    def getDuration(self):
      return self._duration

    def isLoop(self):
        return False

    def getNoteList(self):
        return self._notelist

class NoteList(Expression):
    def __init__(self, note, notelist):
        self._notelist = notelist + [note]

    def getNoteList(self):
      return self._notelist
\end{verbatim}


\newpage

\subsection{parser\_rules.py}
\begin{verbatim}
#!/usr/bin/python
# -*- coding: latin-1 -*-
from parserobjects import *
from lexer_rules import tokens

def p_root(subexpressions):
    'h : tempo compasheader constlistinit voicelist'
    subexpressions[0] = Root(subexpressions[1], subexpressions[2],
     subexpressions[3], subexpressions[4])

def p_root_no_const(subexpressions):
    'h : tempo compasheader voicelist'
    subexpressions[0] = Root(subexpressions[1], subexpressions[2], None,
     subexpressions[3])

def p_tempo(subexpression):
    'tempo : TEMPOBEGIN SHAPE NUM'
    subexpression[0] = Tempo(subexpression[2], int(subexpression[3]))

def p_compasheader(subexpression):
    'compasheader : COMPASHEADERBEGIN NUM SLASH NUM'
    subexpression[0] = CompasHeader(int(subexpression[2]), int(subexpression[4]))

def p_voice(subexpression):
    'voice : VOICEBEGIN LEFTPAR value RIGHTPAR LEFTCURL voicecontent RIGHTCURL'
    subexpression[0] = Voice(subexpression[3], subexpression[6])

def p_compasloop(subexpression):
    'compasloop : LOOPBEGIN LEFTPAR value RIGHTPAR LEFTCURL compaslist
     RIGHTCURL'
    subexpression[0] = CompasLoop(subexpression[3], subexpression[6])

def p_note(subexpression):
    'note : NOTEBEGIN LEFTPAR NOTENAME COMMA value COMMA SHAPE RIGHTPAR
     SEMICOLON'
    subexpression[0] = Note(subexpression[3], None, subexpression[5],
     subexpression[7], False)

def p_note_alter(subexpression):
    'note : NOTEBEGIN LEFTPAR NOTENAME ALTER COMMA value COMMA SHAPE RIGHTPAR
     SEMICOLON'
    subexpression[0] = Note(subexpression[3], subexpression[4],
     subexpression[6], subexpression[8], False)

def p_note_punto(subexpression):
    'note : NOTEBEGIN LEFTPAR NOTENAME COMMA value COMMA SHAPE PUNTO RIGHTPAR
     SEMICOLON'
    subexpression[0] = Note(subexpression[3], None, subexpression[5],
     subexpression[7], True)

def p_note_alter_punto(subexpression):
    'note : NOTEBEGIN LEFTPAR NOTENAME ALTER COMMA value COMMA SHAPE PUNTO
     RIGHTPAR SEMICOLON'
    subexpression[0] = Note(subexpression[3], subexpression[4],
     subexpression[6], subexpression[8], True)

def p_compaslist_base(subexpression):
    'compaslist : compas'
    subexpression[0] = CompasList(subexpression[1], [])

def p_compaslist_rec(subexpression):
    'compaslist : compaslist compas'
    subexpression[0] = CompasList(subexpression[2], subexpression[1].getList())

def p_voice_list_base(subexpression):
    'voicelist : voice'
    subexpression[0] = VoiceList(subexpression[1])

def p_voice_list_rec(subexpressions):
    'voicelist : voicelist voice'
    ### Invierto parametros intencionalmente. voicelist param es opcional en 
    ### el new de la clase
    subexpressions[0] = VoiceList(subexpressions[2], subexpressions[1].getList())

def p_compas(subexpressions):
    'compas : COMPASBEGIN LEFTCURL notelist RIGHTCURL'
    subexpressions[0] = Compas(subexpressions[3].getNoteList())

def p_note_list_base_note(subexpression):
    'notelist : note'
    subexpression[0] = NoteList(subexpression[1], [])

def p_note_list_base_silence(subexpression):
    'notelist : silence'
    subexpression[0] = NoteList(subexpression[1], [])

def p_note_list_rec_note(subexpressions):
    'notelist : notelist note'
    ### Invierto parametros intencionalmente. notelist param es opcional en 
    ### el new de la clase
    subexpressions[0] = NoteList(subexpressions[2], subexpressions[1].getNoteList())

def p_note_list_rec_silence(subexpressions):
    'notelist : notelist silence'
    ### Invierto parametros intencionalmente. notelist param es opcional en 
    ### el new de la clase
    subexpressions[0] = NoteList(subexpressions[2], subexpressions[1].getNoteList())

def p_val_num(subexpression):
    'value : NUM'
    subexpression[0] = int(subexpression[1])


def p_val_cname(subexpression):
    'value : CNAME'
    subexpression[0] = ConstantManager.getInstance().getValue(subexpression[1])
def p_const(subexpressions):
    'const : CONST CNAME EQUALS NUM SEMICOLON'
    subexpressions[0] = Const(subexpressions[2],int(subexpressions[4]), False)

#Una constante que es un puntero a otra constante
def p_const_cname(subexpressions):
    'const : CONST CNAME EQUALS CNAME SEMICOLON'
    subexpressions[0] = Const(subexpressions[2],subexpressions[4], True)

#Es un pasamanos para poder inicializar el ConstantManager y que pueda ser 
#usado por las otras producciones
#(esto asume que las constantes se declaran primero en el header)
def p_const_list_init(subexpressions):
    'constlistinit : constlist'
    subexpressions[0] = subexpressions[1]
    #Sabemos que subexpressions[0] es un constlist, inicializamos el
    # Constantmanager para que el resto
    #de las producciones puedan referenciar constantes.
    #TODO: Pasarle el reserved del Lexer
    ConstantManager.createInstance (subexpressions[0].getList(),[] )


def p_const_list_base(subexpressions):
    'constlist : const'
    subexpressions[0] = ConstList(subexpressions[1],[])

def p_const_list_rec(subexpressions):
    'constlist : constlist const'
    subexpressions[0] = ConstList(subexpressions[2],subexpressions[1].getList())


def p_voice_content_base_loop(subexpressions):
    'voicecontent : compasloop'
    subexpressions[0] = VoiceContent(subexpressions[1],[])

def p_voice_content_base_compas(subexpressions):
    'voicecontent : compas'
    subexpressions[0] = VoiceContent(subexpressions[1],[])

def p_voice_content_rec_compasloop(subexpressions):
    'voicecontent : voicecontent compasloop'
    subexpressions[0] = VoiceContent(subexpressions[2],subexpressions[1].getList())

def p_voice_content_rec_compas(subexpressions):
    'voicecontent : voicecontent compas'
    subexpressions[0] = VoiceContent(subexpressions[2],subexpressions[1].getList())

def p_silence(subexpression):
    'silence : SILENCEBEGIN LEFTPAR SHAPE RIGHTPAR SEMICOLON'
    subexpression[0] = Silence(subexpression[3],None)

def p_silence_punto(subexpression):
    'silence : SILENCEBEGIN LEFTPAR SHAPE PUNTO RIGHTPAR SEMICOLON'
    subexpression[0] = Silence(subexpression[3],True)


def isReserved(token):
    return token in ('const',
    'do',
    're',
    'mi',
    'fa',
    'sol',
    'la',
    'si',
    'tempo',
    'compas',
    'repetir',
    'voz',
    'negra',
    'blanca',
    'redonda',
    'corchea',
    'smicorchea',
    'fusa',
    'semifusa')

def p_error(subexpressions):
    #print ("----------------------------")
    #print ("----------------------------")
    #print (subexpressions)

    if (subexpressions != None):
        if isReserved(subexpressions.value):
            strReservedMsg = '(palabra reservada)'
        else:
            strReservedMsg = ''

        raise Exception("[Parser] Error de sintaxis Linea: {0}, Pos (absoluta):
         {1}, Token: <{2}>{3} ".format(subexpressions.lineno,
          subexpressions.lexpos, subexpressions.value, strReservedMsg))
    else:
        raise Exception("[Parser] Archivo incompleto")
\end{verbatim}
\newpage

\subsection{midi.py}
\begin{verbatim}
#!/usr/bin/python
# -*- coding: latin-1 -*-
from parserobjects import *

class MidiTranslator:
    def generateMIDIFile(self, root, file_output):

        strFormatHeader0 = "MFile 1 {0} 384"
        strFormatHeader1 = "MTrk"
        strFormatHeader2 = "000:00:000 Tempo {0}"
        strFormatHeader3 = "000:00:000 TimeSig {0}/{1} 24 8"
        strFormatHeader4 = "000:00:000 Meta TrkEnd"
        strFormatHeader5 = "TrkEnd"


        strFormatVoyHeader1 = 'MTrk'
        strFormatVoyHeader2 = '000:00:000 Meta TrkName "Voz {0}"'
        strFormatVoyHeader3 = '000:00:000 ProgCh ch={0} prog={1}'


        strFormatVoyHeader4 = '{0:03d}:{1:02d}:{2:03d} Meta TrkEnd'
        strFormatVoyHeader5 = "TrkEnd"

        self.writeMidi(strFormatHeader0.format(len(root.getVoiceList().getList())
                       + 1),file_output)
        self.writeMidi(strFormatHeader1,file_output)
        self.writeMidi(strFormatHeader2.format(self.calculateTempo( 
                       root.getTempo())),file_output)
        self.writeMidi(strFormatHeader3.format(root.getCompasHeader().getNumerator(),
                       root.getCompasHeader().getDenominator()),file_output)
        self.writeMidi(strFormatHeader4,file_output)
        self.writeMidi(strFormatHeader5,file_output)

        vlist = root.getVoiceList().getList()

        for i in range(0, len(vlist)):
            #Encabezado de la Voz
            v = vlist[i]
            channel = i + 1
            self.writeMidi(strFormatVoyHeader1,file_output)
            self.writeMidi(strFormatVoyHeader2.format(channel),file_output)
            self.writeMidi(strFormatVoyHeader3.format(channel,
                           v.getInstrument()),file_output)
            compasId = 0
            vc = v.getCompasses()
            for j in range(0, len(vc)):
                content = vc[j]

                if content.isLoop():
                    repeat = content.getRepeat()
                    contentList = content.getCompasses().getList()
                else:
                    repeat = 1
                    contentList = [content]

                for k in range(0, repeat):
                    for c in contentList:
                        self.writeCompass(c, root, compasId, channel,
                                          file_output)
                        compasId = compasId + 1

            self.writeMidi(strFormatVoyHeader4.format(compasId, 0, 0)
                           ,file_output)

            self.writeMidi(strFormatVoyHeader5,file_output)



    def writeMidi(self, line, file_output):
        file_output.writelines(line)
        file_output.write('\n')

    def writeCompass(self, compass, root, compassId, channel, file_output):

        strFormatVoyNoteOn = '{0:03d}:{1:02d}:{2:03d} On ch={3} note={4}{5}
                              vol=70'
        strFormatVoyNoteOff = '{0:03d}:{1:02d}:{2:03d} Off ch={3} note={4}{5}
                               vol=0'

        clickOffset = 0
        pulseOffset = 0
        nList = compass.getNoteList()
        for i in range (0, len(nList)):
            n = nList[i]
            clicks = int(n.getDuration() *
                         root.getCompasHeader().getDenominator() * 384)
            finalClick = (clicks + clickOffset) % 384

            finalPulse = pulseOffset + (clicks + clickOffset) // 384


            if not n.isSilence():
                noteName = self.fromLatinToAmerican(n.getHeight())
                octave = n.getOctave()
                self.writeMidi(strFormatVoyNoteOn.format(compassId, pulseOffset,
                               clickOffset, channel, noteName, octave), 
                               file_output)

                #Imprimo la nota en el archivo. Si la nota es la ultima del
                #compas, entonces imprimo el Off al comienzo del proximo compas
                if i < len(nList) - 1:

                    self.writeMidi(strFormatVoyNoteOff.format(compassId,
                                   finalPulse, finalClick, channel, noteName,
                                   octave),file_output)
                else:
                    self.writeMidi(strFormatVoyNoteOff.format(compassId + 1, 0,
                                   0, channel, noteName, octave),file_output)


            clickOffset = finalClick
            pulseOffset = finalPulse




    def calculateTempo(self, tempo):
        return int(1000000 * 60 / (tempo.getShapeDuration() * 4 * tempo.getCount()))

    def fromLatinToAmerican(self, noteName):
        #A: la
        #B: si
        #C: do
        #D: re
        #E: mi
        #F: fa
        #G: sol.

        return noteName.replace('la', 'a').replace('si','b').replace('do',
               'c').replace('re', 'd').replace('mi', 'e').replace('fa',
               'f').replace('sol', "g")
\end{verbatim}


\end{document}

